\section{Background and Related Work}
This section discusses earlier research that leveraged code change history, especially at the method level, to analyze and build maintenance models. It then discusses research that focused on building history construction tools.
\subsection{Research on Code Change History and Maintenance}
Software maintenance models commonly target the prediction of changes and bugs, as components that undergo frequent modifications or are tied to bug fixes are often harmful~\cite{chowdhury_good_2025, arvanitou_method_2017, di_grazia_code_2023,opu2025llm}. Early identification of such components can substantially lower future maintenance effort and cost~\cite{gunes_koru_identifying_2007, abbas_software_2020, chowdhury_good_2025, chowdhury_empirical_2022}. Leveraging various code metrics, Koru \textit{et al.}~\cite{gunes_koru_identifying_2007} built tree-based models to forecast change-prone classes. In a similar vein, Romano and Pinzger~\cite{romano_using_2011} applied code metrics to predict change-prone \texttt{Java} interfaces. Khomh \textit{et al.}~\cite{khomh_exploratory_2009} demonstrated that certain types of code smells can also signal classes likely to change. Bug prediction models pursue a similar objective to \textit{change-proneness} prediction models: identifying code elements—such as classes or files—that are more likely to contain bugs in future versions~\cite{aleithan_explainable_2021, alsolai_predicting_2018, basili_validation_1996, gil_correlation_2017, zimmermann_predicting_2007, kamei_defect_2016}. These studies, however, primarily focused on the class or file level, which has been reported as less useful by both practitioners and researchers~\cite{grund_codeshovel_2021, pascarella_performance_2020, giger_method-level_2012, chowdhury_method-level_2024}. For example, finding bugs in a class or file can be difficult due to its large size. Consequently, method-level change and bug prediction models have gained traction in recent years.

Method-level bug prediction has been studied by Giger \textit{et al.}~\cite{giger_method-level_2012}, Mo \textit{et al.}~\cite{mo_exploratory_2022}, Ferenc \textit{et al.}~\cite{ferenc_automatically_2020}, Shippey \textit{et al.}~\cite{shippey_so_2016}, Hata \textit{et al.}~\cite{hata_bug_2012}, Menzies \textit{et al.}~\cite{menzies_data_2007}, Pascarella \textit{et al.}~\cite{pascarella_performance_2020}, and Chowdhury \textit{et al.}~\cite{chowdhury_method-level_2024}. These studies heavily relied on constructing method-level change histories—either to build labeled datasets or to derive change metrics used as input features for machine learning models. For instance, Mo \textit{et al.}~\cite{mo_exploratory_2022} employed metrics such as the number of added and deleted lines within a method to predict its bug-proneness, a strategy also adopted by Pascarella \textit{et al.}~\cite{pascarella_performance_2020}. Chowdhury \textit{et al.}~\cite{chowdhury_method-level_2024} used method-level change history to identify bug-fixing commits and label methods as buggy or non-buggy. The accuracy of the tools used to construct method-level change histories is therefore critical, as inaccuracies can lead to erroneous change metrics and incorrect bug labels.

\subsection{Research on History Construction}
Earlier studies on method-level tracking in \texttt{Git} used separate files to record method changes alongside \texttt{Git’s} file-level history. After extensive preprocessing of the entire codebase, this approach enables the use of \texttt{git log} to retrieve method histories. Hata \textit{et al.}~\cite{hata_historage_2011} introduced this technique with \textit{Historage}, which tracks changes to constructors, methods, and fields. However, Higo \textit{et al.}~\cite{higo_tracking_2020} found that \textit{Historage} struggles to track small methods when renamed or moved across files. To overcome these issues, they proposed \textit{FinerGit}, which represents methods with one token per line and applies heuristics for change comparison. Still, the preprocessing is time-intensive and incurs significant storage overhead—an issue common across similar tools~\cite{kim_when_2005,tu_integrated_2002,zimmermann_fine-grained_2006,hassan2004c}. Other tools~\cite{huang_cldiff_2018,falleri_fine-grained_2024} analyze changes between two versions but fall short of supporting origin analysis required by developers~\cite{grund_codeshovel_2021} and researchers~\cite{chowdhury_good_2025}.

Recent studies have proposed solutions that eliminate the need for upfront repository preprocessing and avoid tracking methods at each \texttt{Git} commit~\cite{jodavi_accurate_2022, grund_codeshovel_2021}. Grund \textit{et al.}~\cite{grund_codeshovel_2021} introduced \textit{CodeShovel}, a state-of-the-art tool capable of recovering method histories despite various transformations. On a manually annotated oracle, \textit{CodeShovel} outperformed both research tools (e.g., \textit{FinerGit}) and industry tools (e.g., \textit{IntelliJ}, \textit{GitLog}). However, Jodavi \textit{et al.}~\cite{jodavi_accurate_2022} found that \textit{CodeShovel} struggles when methods undergo significant body changes, such as during moves, often failing to track them correctly or misattributing their origin.

By leveraging the RefactoringMiner~\cite{tsantalis_refactoringminer_2022}, the authors also noted inaccuracies in the original \textit{CodeShovel} oracle and revised it based on their interpretation to create a more reliable benchmark. In general, they identified inaccuracies in 27 method histories due to method extraction refactorings. The authors updated the \textit{CodeShovel} oracle based on RefactoringMiner’s suggestions, but this approach has two key issues: (i) the accuracy of the remaining 173 methods is unknown, and (ii) RefactoringMiner itself has low recall (54.2) for the extract and move method~\cite{tsantalis_refactoringminer_2022}, raising doubts about the updates’ reliability. The authors also developed \textit{CodeTracker}, which employs RefactoringMiner’s advanced entity matching algorithm~\cite{tsantalis_refactoringminer_2022}. When evaluated with the new oracle, \textit{CodeTracker} showed significant improvements in precision and recall over \textit{CodeShovel}, but its reliance on RefactoringMiner introduced a performance bottleneck, making it significantly slower.

\emph{\textbf{Motivation:}} Like its predecessor, the new \textit{CodeTracker} oracle remains subject to internal validity threats, as it relies on the RefactoringMiner and on researchers’ interpretations. A more systematic oracle construction process is needed for the reliable evaluation of method history tools. Additionally, there is a need for a new algorithm that maintains high accuracy without sacrificing runtime performance, unlike \textit{CodeTracker}.
